\section{Experimental Model}

The design separates three effects that are otherwise confounded:

\begin{enumerate}
\item the presence of an auditory stimulus;
\item rhythmic or attentional pacing;
\item left-right alternation.
\end{enumerate}

Three matched conditions are therefore required:

\begin{description}
\item[A: alternating bilateral.] The same auditory event alternates between left and right channels according to a fixed temporal pattern.
\item[B: matched non-alternating.] The same event timing, spectrum, loudness, duration, and total acoustic energy are delivered without systematic left-right alternation.
\item[C: quiet regulation.] The same regulation instructions are delivered without the experimental stimulus.
\end{description}

The primary specific bilateral contrast is \(A-B\). The contrast \(B-C\) estimates a nonspecific auditory/rhythmic effect. The contrast \(A-C\) is descriptive unless the component-specific contrast \(A-B\) is also supported.

For eligible regulation event \(e\), let \(T_{\mathrm{rec}}^{(e)}\) be the core ARC-ADHD monitoring-anchored recovery latency. The auditory conditions must have matched exposure duration so that any condition difference in the pre-monitoring interval \(T_{\mathrm{pre}}\) is not produced trivially by unequal stimulus length. \(T_{\mathrm{pre}}\) and total return time \(T_{\mathrm{total}}\) are reported separately. The bilateral-specific hypothesis predicts a lower participant-weighted restricted mean recovery burden for A than for B over the same monitoring horizon, as defined in the analysis section. This prediction also requires acceptable tolerability; a shorter latency among only the children who complete sound exposure is insufficient. Faster recovery in both A and B relative to C, without a difference between A and B, supports a general auditory or pacing effect rather than a bilateral-specific mechanism.
